%% kapa: exam=1 round=28 year=2017 period=1 form=A subject=economics subject_ko=경제학원론 seq=10 number=50
%% source: Q-Net 감정평가사 기출문제 (cst00309) artlSeq=5199603 file=extracted/1008/5199603/2017년 제28회 감정평가사 1차시험 문제지/1교시 A형.pdf
%% conversion: kordoc 4.24.51 -> tools/gen1.py (auto); review=proofread-sonnet; flags=none
\begin{question}{50}
어느 마을에 주민들이 염소를 방목할 수 있는 공동의 목초지가 있다. 염소를 방목하여 기를 때 얻는 총수입은 $R=10(20X-X^{2})$이고, 염소 한 마리에 소요되는 비용은 20이다. 만약 개별 주민들이 아무런 제한 없이 각자 염소를 목초지에 방목하면 마을 주민들은 총 $X_{1}$ 마리를, 마을 주민들이 마을 전체의 이윤을 극대화하고자 한다면 총 $X_{2}$ 마리를 방목할 것이다. $X_{1}$과 $X_{2}$는? (단, $X$는 염소의 마리수이다.)
\begin{choices}
  \item 12, 9
  \item 12, 16
  \item 16, 12
  \item 18, 9
  \item 18, 12
\end{choices}
\end{question}
\answer{4}
